\documentclass[11pt,letterpaper]{article}
\usepackage[T1]{fontenc}
\usepackage[utf8]{inputenc}
\usepackage[scaled=0.95]{helvet}
\renewcommand{\familydefault}{\sfdefault}
\usepackage[top=1.7cm,bottom=1.5cm,left=2.2cm,right=2.2cm]{geometry}
\usepackage{microtype}
\usepackage[hidelinks]{hyperref}
\pagestyle{empty}
\setlength{\parindent}{0pt}
\setlength{\parskip}{7pt}

\begin{document}

{\large\bfseries Grover Aruquipa}\par\vspace{-6pt}
{\small Qu\'ebec, QC, Canada \ \textbullet\ \ grover.aruquipa.9@gmail.com \ \textbullet\ \ groveraruquipa.github.io}

\vspace{4pt}
Robotics and Grasping Team, Apera AI, Vancouver\\
Mechatronics Engineer, Robotics and Grasping

Dear Robotics and Grasping Team,

I am applying for the Mechatronics Engineer position on your Robotics and Grasping team. My bachelor thesis was a reconfigurable Delta robot for product picking, I worked as a computer vision engineer on industrial anomaly detection, and for the last three years I have designed and built robots with integrated grippers. A role that brings machine vision, mechanical design and grasping together in one product is the one I am looking for.

I am finishing my PhD in Mechanical Engineering at Universit\'e Laval with Prof.\ Cl\'ement Gosselin and Prof.\ Patrice Lambert, and I defend in January 2027. I designed, built and controlled three parallel manipulators with 4, 7 and 8 actuators, with every motor on the base and the gripper driven through the mechanism. I took each one from the mathematical model and the SolidWorks CAD through fabrication to 500\,Hz closed-loop control and experimental validation. My 4-DoF coaxial wrist-gripper, with 0.25\,kg of moving mass in a 180\,mm footprint, completed ten continuous tool revolutions with sub-degree tracking.

My work combines the mechanical and the software side of a robot. I wrote real-time control software in C++ during my master's at FEMTO-ST, I work in Python, MATLAB and Simulink, I have used OpenCV, ROS and MoveIt, and I validate every kinematic model and trajectory in MuJoCo before running it on the real robot. I am now building a Delta robot with an integrated Sarrus gripper, aimed at sensing the grasp force through the motor currents to handle delicate objects.

I also enjoy working with others. I led two graduate students from concept to working prototype on an 8-actuator robot, and I led the Universit\'e Laval team that won first place in the graduate division of the ASME Student Mechanisms and Robotics Design Competition in 2025.

I will be eligible for an open Post-Graduation Work Permit as soon as I complete my PhD, so I do not need sponsorship, and I am ready to relocate to Vancouver and work in person from January 2027. Prof.\ Gosselin, Prof.\ Lambert and Univ.-Prof.\ Andreas M\"uller of JKU Linz can speak to my work; their details are in my CV.

Thank you for your time. I would be glad to show you my robots running, on video or live.

Sincerely,\\
Grover Aruquipa

\end{document}
